% Standalone source-record documentation; no research attempts.
\documentclass[11pt]{article}
\usepackage[margin=25mm]{geometry}
\usepackage{fontspec,unicode-math}
\setmainfont{Latin Modern Roman}
\setmathfont{Latin Modern Math}
\usepackage{amsmath,amssymb,mathtools,xurl,hyperref}
\hypersetup{colorlinks=true,urlcolor=blue,linkcolor=blue}
\setcounter{secnumdepth}{0}
\setlength{\parindent}{0pt}
\setlength{\parskip}{5pt}
\setlength{\emergencystretch}{3em}
\newcommand{\sref}[2]{\hyperlink{source:#1:#2}{[S#2]}}
\newcommand{\cataloguescope}[1]{\paragraph{Source record.}#1}
\newcommand{\floor}[1]{\left\lfloor#1\right\rfloor}
\begin{document}
% UnsolvedMath ID 30004925; source catalogue code OWR-8415359-015
\setcounter{section}{30004924}
\section{Size-Efficient Circuits Resilient to Short-Circuit Errors}\label{30004925}
{\small\sffamily UnsolvedMath ID 30004925 · Computer Science}
\subsection{Definitions and mathematical statement}

\paragraph{Shared notation.}Standard number systems are $\mathbb N,\mathbb Z,\mathbb Q,\mathbb R,\mathbb C$. Specialized notation and index ranges follow the cited source and the definitions below.


In a short-circuit fault, a Boolean gate outputs one of its input values instead of its intended output. A fault-tolerant construction must specify which gates or paths the adversary may corrupt and how circuit-size overhead is measured.
\paragraph{Record scope and qualifications.}
The catalogue title is retained as metadata. Named certificates, contracts, bounds or frameworks in that title are not asserted as results of the original source extract.
The primary report already gives a quasipolynomial-size construction under a bounded-fraction path-fault model; the cleaned record does not state a definite remaining overhead improvement or its fault threshold. It is retained as a source formulation, not asserted as a fresh unresolved existence claim.

Source-record limitation: The primary source already proves a quasipolynomial-size resilience result; the cleaned question does not state a definite remaining size improvement or fault threshold.
\paragraph{Statement recorded in the supplied source.}
Given a Boolean circuit $C$, can one construct a circuit $C'$ of comparable size that computes the same function even when an adversary replaces the output of some gates by one of their inputs? \sref{30004925}{2}
\paragraph{Source status.}
The supplied catalogue labels this source record partially solved; that metadata was checked against the supplied snapshot on 2026-10-11. The label is preserved for provenance and is not an independent verification of present open status or mathematical correctness. The catalogue literature assessment is dated 2026-08-22.
\subsection{Short English statement}
Ask for comparable-size protection of Boolean computations against the specified fault type.
\subsection{Sources}
\begin{itemize}
\item[\hypertarget{source:30004925:1}{S1}] Ulam AI and the UnsolvedMath Contributors, UnsolvedMath: A Curated Collection of Open Mathematics Problems, record 30004925 (OWR-8415359-015). Catalogue curation and metadata: CC BY 4.0. The source extract is retained; definitions, English summaries and qualifications are editorial additions.
\url{https://huggingface.co/datasets/ulamai/UnsolvedMath}.
\item[\hypertarget{source:30004925:2}{S2}] Complexity Theory (2022), Oberwolfach Reports, DOI 10.4171/owr/2021/54, pp. 2955--3005
\url{https://doi.org/10.4171/owr/2021/54}.
\item[\hypertarget{source:30004925:3}{S3}] Efremenko et al., Circuits Resilient to Short-Circuit Errors
\url{https://par.nsf.gov/servlets/purl/10359538}.
\item[\hypertarget{source:30004925:4}{S4}] STOC 2022 publication
\url{https://doi.org/10.1145/3519935.3520007}.
\item Source-specific Oberwolfach Reports licences, including CC BY-SA 4.0 where applicable, are retained; see THIRD\_PARTY\_NOTICES.md.
\end{itemize}
\paragraph{Preparation and review.}
The original source and UnsolvedMath contributors retain attribution for the source extract. Codex (GPT-6; exact serving revision not exposed) assisted with definitions, English summaries, scope notes and formatting. No new human mathematical authorship or independently verified solution is asserted. Permanent source ID 30004925 and code OWR-8415359-015 checked against the complete supplied registry and current repository. Distinct upstream records, including repeated questions and broad or damaged extracts, are retained for complete coverage of the source CS category. A repeated record is not a new independent conjecture. Complete mathematical-equivalence and current-literature review has not been performed. Preparation and final review follow contribution prompts C and D, with source gaps disclosed. Structural and rendering checks do not establish a new theorem or open status.
\label{end:30004925}
\end{document}
